电与磁最初分别由静电学、稳恒电流和电磁感应来描述。图~\ref{review-maxwell-portrait} 中的麦克斯韦（James Clerk Maxwell）的关键工作，是看出这些规律必须组成
一个彼此相容的整体；而卡末林·昂内斯（Heike Kamerlingh Onnes）在低温实验中发现的超导现象，则把电磁场与物质的相互作用推到了一个
极端而清晰的情形：电阻可以消失，磁通却不能只按“零电阻导体”来理解。图~\ref{review-history-context} 将这两条历史线索并列，左侧是场的理论，右侧是场与低温物质的实验。

电磁感应定律揭示了变化的磁场能够产生电场。由此自然会追问：变化的电场能否产生磁场？若原有的
Ampere 定律只能处理稳恒电流，它在含有随时间变化的电荷与电场时是否仍然完备？Maxwell 在尚无直接
实验依据的情况下，从理论自洽性出发补上位移电流项，使电场和磁场能够彼此激发。由此得到的电磁波
传播速度与当时测得的光速一致，他进一步提出光的电磁理论，经典电磁学的基本框架遂告建立。

这一过程值得特别注意：Maxwell 并不是把彼此无关的四条实验定律简单抄在一起，而是要求它们同时
满足电荷守恒、空间对称性和有限传播速度。静电场的 Gauss 定律告诉我们电荷能够成为电场的源，磁场
Gauss 定律则表明磁感线没有已经被发现的起点和终点；Faraday 定律说明变化磁场会在周围激起涡旋电场。
唯独旧的 Ampere 定律只容纳稳恒传导电流，在时间变化问题中显得不对称。位移电流修正消除了这种
不对称，并使方程组能够在没有电荷、没有导线的真空中仍然具有随时间传播的非零解。电磁波并非额外
塞进理论的假设，而是完备方程组自身推出来的结果。

从学习方法上说，本讲的难点不在记住更多公式，而在分清三种层次。第一层是实验事实，例如电荷产生
电场、电流产生磁场以及电磁感应；第二层是把事实压缩为局部方程，并检查这些方程能否同时成立；第三层
才是从方程推出波速、能流和辐射。只记“变化电场产生磁场”会漏掉为什么必须有这一项；只记四条方程
又会看不出电磁波从何而来。后文的每一步都围绕“源—场—传播—能量”这条线索展开。

\maxwellOnnesContext

\historyfigure[0.27\linewidth]{maxwell-built-in-v1.png}%
{图 10.1：麦克斯韦（James Clerk Maxwell，1831--1879），历史人物肖像。\label{review-maxwell-portrait}}

\subsection{位移电流：补上安培定律的缺口}

本段限定真空或磁导率近似为 $\mu_0$ 的非磁介质，$\vect j$ 为自由传导电流密度，$\rho$ 为自由电荷密度。含磁化介质时应使用 $\vect H$，不能仍把任意介质中的 $\vect B$ 直接除以 $\mu_0$。
稳恒电流产生磁场时，安培环路定律为
\begin{equation}
  \oint_C\vect B\cdot\dd\vect l=\mu_0 I_{\rm c},
\end{equation}
其中 $I_{\rm c}=\int_S\vect j\cdot\dd\vect S$ 是穿过跨越面 $S$ 的传导电流。若电流处处稳恒，选择
同一边界 $C$ 的不同跨越面，所得 $I_{\rm c}$ 相同，公式没有歧义。

先看第一个矛盾。设一块均匀、弱导电的球形介质中心原有电荷，电荷随时间沿径向向外泄漏。
电流密度具有严格的球对称性。若在球内任取一条闭合环路，按原来的 Ampere 环路定律，只要跨越面上
有净电流穿过，环路周围似乎就应当出现环向磁场。然而球对称系统中不存在一个在每一点都非零、又能
选出固定切向的矢量场；任何非零环向磁场都会破坏球对称性。因此，单独用传导电流作磁场源会导出
相互冲突的结论。矛盾的根源是：电荷流出时，球内电荷量和由 Gauss 定律决定的电场也同时在变化。

把这个论证写得更定量一些。若半径 $r$ 的球面内电荷为 $Q_r(t)$，球对称性给出
\begin{equation}
  4\pi r^2D_r=Q_r(t).
\end{equation}
设向外传导电流为 $I_r$，电荷守恒要求 $I_r=-\dd Q_r/\dd t$；另一方面，穿过同一球面的位移电流为
\begin{equation}
  I_{\rm d}=\frac{\dd}{\dd t}(4\pi r^2D_r)=\frac{\dd Q_r}{\dd t}=-I_r.
\end{equation}
球对称性还把这个积分结论加强为每一点的关系：
\[
 j_r=\frac{I_r}{4\pi r^2},\qquad
 \frac{\partial D_r}{\partial t}=-\frac{I_r}{4\pi r^2},
 \qquad \vect j+\frac{\partial\vect D}{\partial t}=0.
\]
因此不是只在闭合面上“总量抵消”，而是局部的两种源也相消。
无外加磁场时，对称性只允许径向磁场；再由任意同心球面磁通为零，可得 $B_r=0$。于是球对称解确实为 $\vect B=0$。

图 10.2 中实线箭头表示径向电场，灰色虚线箭头表示径向传导电流。任选闭合环路 $\Gamma$ 都不能在
球对称条件下挑出唯一的环向磁场；图右把同一点的 $\vect j$ 与 $\partial\vect D/\partial t$ 单独放大，两者大小相等、方向相反。

\radialCurrentSphereDiagram

给电容器充电时，问题暴露出来：取 $S_1$ 穿过导线，有 $I_{\rm c}=I$；取 $S_2$ 从两极板之间鼓出，
它不切断任何导线，因而 $I_{\rm c}=0$。两个面有同一边界，却会给出两个不同的磁场环量。

这里真正参与定律的是“同一闭合回路的任意跨越面”。环路 $C$ 可以想成一根铁丝围成的圈，$S_1$ 与
$S_2$ 则像两张边缘都粘在铁丝上的橡皮膜。一张膜穿过导线，另一张膜向电容器间隙鼓起；二者形状不同，
边界却完全相同。若右端只计算传导电流，左端同一个 $\oint_C\vect B\cdot\dd\vect l$ 就会得到两个值，
这不是近似误差，而是定律内部失去定义。位移电流项的任务正是让右端也只依赖边界，而不依赖所选曲面。

图~\ref{review-surface-choice} 中两种蓝色剖面具有相同边界；从左图变为右图时，跨越面绕开左极板而不切断导线。

\displacementCurrentDiagram

极板间虽然没有自由电荷穿越，电位移通量却在改变。Maxwell 定义位移电流
\begin{equation}
  I_{\rm d}=\frac{\dd}{\dd t}\int_S\vect D\cdot\dd\vect S,
\end{equation}
一般介质中的形式为 $\oint_C\vect H\cdot\dd\vect l=\int_S(\vect j+\partial\vect D/\partial t)\cdot\dd\vect S$。在本段的非磁介质条件下，它写为
\begin{equation}
  \boxed{\oint_C\vect B\cdot\dd\vect l
  =\mu_0\left(\int_S\vect j\cdot\dd\vect S
  +\frac{\dd}{\dd t}\int_S\vect D\cdot\dd\vect S\right)}.
\end{equation}
真空中 $\vect D=\varepsilon_0\vect E$。平行板电容器忽略边缘效应时，板间
$E=Q/(\varepsilon_0 A)$，于是
\begin{equation}
  I_{\rm d}=\varepsilon_0\frac{\dd(EA)}{\dd t}
  =\frac{\dd Q}{\dd t}=I.
\end{equation}
因此 $S_1$ 上的传导电流与 $S_2$ 上的位移电流恰好相等，磁场环量不再依赖跨越面的选取。

若在极板之间取半径为 $r$、与极板同轴的圆形环路，并假设 $r$ 小于极板半径，则该环路所围面积中的
位移电流为 $I_{{\rm d},r}=I\pi r^2/A$。由圆柱对称性，环路上的磁感应强度大小处处相同且沿切向，故
\begin{equation}
  B(2\pi r)=\mu_0I\frac{\pi r^2}{A},
  \qquad B=\frac{\mu_0Ir}{2A}.
\end{equation}
在极板投影之外，所包围的总位移电流已等于 $I$，于是 $B=\mu_0I/(2\pi r)$。这表明充电间隙中的
磁场并非抽象的“修补量”：它同样可由环路定律给出具体空间分布，并与导线周围磁场连续衔接。

回看弱导电球的例子，穿过闭合面的向外传导电流使内部电荷减少，与此同时电通量也减少。传导电流与
位移电流的代数和满足连续性，因而不会再凭空选出一个破坏球对称性的磁场方向。两个表面上看不同的
矛盾，其共同点正是都遗漏了变化电场的磁效应。

这个修正并非只为解决电容器的特殊矛盾。局部电荷守恒要求
\begin{equation}
  \nabla\cdot\vect j=-\frac{\partial\rho}{\partial t}.
\end{equation}
若仍用 $\nabla\times\vect B=\mu_0\vect j$，两边取散度便得到
$0=\mu_0\nabla\cdot\vect j$，它只允许电荷密度不随时间变化。加入位移电流后，利用
$\nabla\cdot\vect D=\rho$，有
\begin{equation}
  0=\mu_0\left(\nabla\cdot\vect j+
  \frac{\partial}{\partial t}\nabla\cdot\vect D\right)
  =\mu_0\left(\nabla\cdot\vect j+\frac{\partial\rho}{\partial t}\right),
\end{equation}
恰好还原连续性方程。

图~\ref{review-continuity} 从上向下补入位移电流后，散度方程恢复一般的电荷守恒。

\continuityDiagram

这里的“位移电流”不是极板间真的有电荷横渡真空；它是变化电场对磁场的贡献，单位与电流相同。
这个名称带有明显的历史色彩。Maxwell 所处的时代仍常把电磁场设想为某种介质的状态，认为外场会使
假想的真空介质或真实电介质发生极化，极化状态的变化仿佛形成一种电流。今天已经不再需要这种“以太”
图景，但“位移电流”这一名称以及电位移矢量 $\vect D$ 被保留下来。现代观点强调的不是机械介质真的
流动，而是 $\partial\vect D/\partial t$ 与传导电流一样，是涡旋磁场的源。

位移电流与传导电流还应在物理机制上区分。金属导线中的传导电流伴随自由载流子的定向运动，可以在
截面两侧统计净通过的电荷；真空电容间隙中的位移电流则没有这样的电荷输运。二者之所以能写在同一个
括号中，是因为它们对磁场环量产生同样的源效应，并共同保证总电流连续。不能据此断言“真空中也有
电子从一块极板跳到另一块极板”，也不能把位移电流当作某种额外漏电。


% EM teaching enrichment 20260927
\paragraph{位移电流不是电子穿过真空间隙}
平行板充电时，导线中传导电流改变极板自由电荷，板间的变化电场对应位移电流。二者在磁场环量关系中以相同单位出现，却描述不同机制。
对于均匀线性介质，适用的宏观项为 $\partial\boldsymbol D/\partial t$，包括真空项与极化变化；若介质还漏电，则传导项也要保留。
检查一个跨越面时应分别问：有没有自由电荷穿面？穿面的电位移通量是否变化？只凭“这里没有导线”不能把安培--麦克斯韦方程的右边判为零。

\subsection{Maxwell 方程组}

在真空的宏观描述中，电磁场的四条基本方程可写成
\begin{align}
  \nabla\cdot\vect E&=\frac{\rho}{\varepsilon_0},
  &\oint_{\partial V}\vect E\cdot\dd\vect A&=\frac{Q_V}{\varepsilon_0},
  \\
  \nabla\cdot\vect B&=0,
  &\oint_{\partial V}\vect B\cdot\dd\vect A&=0,
  \\
  \nabla\times\vect E&=-\frac{\partial\vect B}{\partial t},
  &\oint_C\vect E\cdot\dd\vect l&=-\frac{\dd}{\dd t}\int_S\vect B\cdot\dd\vect S,
  \\
  \nabla\times\vect B&=\mu_0\vect j+\mu_0\varepsilon_0
  \frac{\partial\vect E}{\partial t},
  &\oint_C\vect B\cdot\dd\vect l&=\mu_0 I_{\rm c}+\mu_0\varepsilon_0
  \frac{\dd}{\dd t}\int_S\vect E\cdot\dd\vect S.
\end{align}
左列是微分形式，描述一个点附近的场；右列是积分形式，直接联系闭合曲面或闭合回路上的整体量。
四条方程分别表达：电荷是电场的源；没有观测到孤立磁荷；变化的磁场产生涡旋电场；传导电流与
变化的电场共同产生涡旋磁场。

积分形式适合直接分析具有对称性的有限区域，微分形式则把定律落实到空间中的每一点。例如
$\nabla\cdot\vect E=\rho/\varepsilon_0$ 在直角坐标系中就是
\begin{equation}
  \frac{\partial E_x}{\partial x}+\frac{\partial E_y}{\partial y}
  +\frac{\partial E_z}{\partial z}=\frac{\rho}{\varepsilon_0},
\end{equation}
而 Faraday 定律包含三个按 $x\to y\to z\to x$ 循环轮换得到的分量式，例如
\begin{equation}
  \frac{\partial E_z}{\partial y}-\frac{\partial E_y}{\partial z}
  =-\frac{\partial B_x}{\partial t}.
\end{equation}
安培--Maxwell 定律同样给出三个分量关系，例如
\begin{equation}
  \frac{\partial B_z}{\partial y}-\frac{\partial B_y}{\partial z}
  =\mu_0j_x+\mu_0\varepsilon_0\frac{\partial E_x}{\partial t}.
\end{equation}
这些式子可由积分定律分别作用在无限小的 Gauss 盒或矩形回路上得到。中学竞赛很少要求逐分量求解
完整方程组，但看清分量间的循环结构，有助于理解后面的平面波推导不是另立的新规律。

完整写出两个旋度方程的分量结构，可以看到这种循环关系：
\begin{align}
 \frac{\partial E_z}{\partial y}-\frac{\partial E_y}{\partial z}
   &=-\frac{\partial B_x}{\partial t},\\
 \frac{\partial E_x}{\partial z}-\frac{\partial E_z}{\partial x}
   &=-\frac{\partial B_y}{\partial t},\\
 \frac{\partial E_y}{\partial x}-\frac{\partial E_x}{\partial y}
   &=-\frac{\partial B_z}{\partial t},\\
 \frac{\partial B_z}{\partial y}-\frac{\partial B_y}{\partial z}
   &=\mu_0j_x+\mu_0\varepsilon_0\frac{\partial E_x}{\partial t},\\
 \frac{\partial B_x}{\partial z}-\frac{\partial B_z}{\partial x}
   &=\mu_0j_y+\mu_0\varepsilon_0\frac{\partial E_y}{\partial t},\\
 \frac{\partial B_y}{\partial x}-\frac{\partial B_x}{\partial y}
   &=\mu_0j_z+\mu_0\varepsilon_0\frac{\partial E_z}{\partial t}.
\end{align}
每一个旋度分量都描述“垂直于某一小面的环流密度”。这也解释了为什么平面波中只需选择一对互相
垂直的电、磁分量，另外几条分量式就自动化为 $0=0$。

方程组还需要 Lorentz 力定律
\begin{equation}
  \boxed{\vect F=q(\vect E+\vect v\times\vect B)}
\end{equation}
才能说明场如何作用于带电粒子。Maxwell 方程决定“场怎样产生和传播”，Lorentz 力决定“场怎样推动
物质”，二者分工不同。

还应避免把“场的方程”和“粒子的运动方程”混为一谈。给定 $\rho,\vect j$ 与边界条件，Maxwell 方程
求出 $\vect E,\vect B$；再把所得场代入 Lorentz 力，结合 Newton 定律，才能求粒子的轨迹。反过来，
粒子运动又会改变电荷、电流分布并重新产生场。完整的电磁动力学是这种相互耦合，而不是先验地把场
当成不会受带电粒子影响的固定背景。

图~\ref{review-maxwell-light} 中连接三个方框的箭头分别表示联立求解和比较传播速度，不是对实验定律的简单罗列。

\maxwellLightPanel


% EM teaching enrichment 20260927
\paragraph{四条方程与材料关系各管什么}
电的通量方程约束电荷源，磁的通量方程排除磁单极源；两条旋度方程描述变化场及电流产生的环流。它们需与电荷运动方程、初值、边界条件共同使用，不是四个互不相关的算式。
例如对安培--麦克斯韦方程取散度，左侧为零，代入电的高斯定律后得到
$\partial\rho/\partial t+\nabla\cdot\boldsymbol J=0$。若删去位移电流项，就会错误地强迫任何问题都满足 $\nabla\cdot\boldsymbol J=0$，无法描述极板正在积累电荷的过程。
介质中还需给出 $\boldsymbol D(\boldsymbol E)$、$\boldsymbol B(\boldsymbol H)$ 等材料关系；这些关系不是所有材料都相同的基本场方程。

\subsection{真空中的电磁波}

在没有电荷和传导电流的真空区域，$\rho=0,\vect j=0$。设一列平面波沿 $+x$ 方向传播，场具有共同相位
\begin{equation}
  \vect E=\vect E_0\cos(kx-\omega t),
  \qquad
  \vect B=\vect B_0\cos(kx-\omega t).
\end{equation}
由 $\nabla\cdot\vect E=0$ 得 $-kE_{0x}\sin(kx-\omega t)=0$，故 $E_{0x}=0$；同理由
$\nabla\cdot\vect B=0$ 得 $B_{0x}=0$。因此源自由空间中的平面电磁波是横波。

不妨选 $\vect E=E_0\cos(kx-\omega t)\uvec y$，并设
$\vect B=B_0\cos(kx-\omega t)\uvec z$。Faraday 定律的 $z$ 分量给出
\begin{equation}
  \frac{\partial E_y}{\partial x}=-\frac{\partial B_z}{\partial t}
  \quad\Longrightarrow\quad
  kE_0=\omega B_0.
\end{equation}
安培--Maxwell 定律的 $y$ 分量为
\begin{equation}
  -\frac{\partial B_z}{\partial x}
  =\mu_0\varepsilon_0\frac{\partial E_y}{\partial t}
  \quad\Longrightarrow\quad
  kB_0=\mu_0\varepsilon_0\omega E_0.
\end{equation}
消去振幅便得到
\begin{equation}
  \boxed{\frac{\omega}{k}=\frac{1}{\sqrt{\mu_0\varepsilon_0}}\equiv c},
  \qquad
  \boxed{E_0=cB_0}.
\end{equation}
这个由电磁常量预言的速度等于真空光速，是“光就是电磁波”的核心证据。更一般地，三者满足
\begin{equation}
  \boxed{\vect B=\frac{1}{\omega}\vect k\times\vect E},
  \qquad
  \vect E\perp\vect B,\quad \vect E\perp\vect k,
  \quad \vect B\perp\vect k.
\end{equation}

推导中场的振幅 $E_0,B_0$ 在方程两边被约去，这并不意味着振幅为零，而是说明真空常量只决定传播
速度及两种场的振幅比，实际振幅由波源强弱决定。频率同样由波源的振荡频率决定；一旦频率 $f$ 给定，
真空中的波长便不能任意选择，而必须满足
\begin{equation}
  \lambda=\frac{c}{f},\qquad k=\frac{2\pi}{\lambda},\qquad \omega=2\pi f.
\end{equation}
传播方向的垂直平面有两个彼此独立的偏振方向。任意线偏振可以分解成这两个方向的叠加；若两分量有
相位差，则可形成椭圆偏振或圆偏振。对这里的真空平面波，$\vect E$ 与 $\vect B$ 同相，三者按
$\vect E,\vect B,\vect k$ 的次序构成右手系。

也可以不先猜传播速度，而直接从 Maxwell 方程消去磁场。对 Faraday 定律取旋度，使用恒等式
$\nabla\times(\nabla\times\vect E)=\nabla(\nabla\cdot\vect E)-\nabla^2\vect E$；在真空无源区
$\nabla\cdot\vect E=0$，再代入安培--Maxwell 定律，得到
\begin{equation}
  \nabla^2\vect E-\mu_0\varepsilon_0\frac{\partial^2\vect E}{\partial t^2}=0.
\end{equation}
磁场同理满足完全相同的波动方程。与标准形式
$\nabla^2\psi-(1/v^2)\partial^2\psi/\partial t^2=0$ 对照，立即读出
$v=1/\sqrt{\mu_0\varepsilon_0}$。这一推导说明速度结论并不依赖于选取沿 $x$ 传播、沿 $y$ 偏振的
特殊坐标；特殊坐标只是把矢量方程变成容易计算的分量式。

把平面简谐波代回波动方程可得 $\omega^2=c^2k^2$。正频率支对应沿 $+\vect k$ 方向传播，改变相位
符号则得到反向传播的波。波动方程是线性的：若两列波分别是解，它们的和仍是解，因此干涉、驻波和
偏振叠加都自然成立。所谓“振幅任意”只是在给定线性真空方程时，方程本身不替波源规定振幅；实际
系统中的天线功率、边界条件和能量守恒仍会决定可实现的场强。类似地，频率可由波源选择，但波速固定，
于是不同频率对应不同波长。可见光只是电磁波谱的一小段，改变频率并不改变其电磁波本性。

图~\ref{review-wave-3d} 中电场位于蓝色平面，磁场位于灰色平面；其方向在下文逐项说明。

\planeWaveThreeDDiagram

图中蓝色曲线位于 $x$--$y$ 平面，灰色曲线位于 $x$--$z$ 平面；同一 $x$ 处的两根场矢量同相。
右手从 $\vect E$ 旋向 $\vect B$，拇指所指就是传播方向。图上的斜向 $z$ 轴只是纸面投影，不表示
磁场真的“倾斜”到传播方向。若 $\vect E$ 始终沿固定横向振动，这列波为线偏振波；若两个横向分量
存在相位差，电场矢量端点可描出椭圆或圆。

\subsection{场的能量与 Poynting 矢量}

真空中电场与磁场的能量密度分别为
\begin{equation}
  u_E=\frac12\varepsilon_0E^2,
  \qquad
  u_B=\frac{B^2}{2\mu_0},
  \qquad u=u_E+u_B.
\end{equation}
为了找出能量怎样流动，将 $\vect E$ 点乘安培--Maxwell 方程、将 $\vect B/\mu_0$ 点乘 Faraday 方程，
除去公共常数后，两个点乘式分别为
\begin{align}
 \frac1{\mu_0}\vect E\cdot(\nabla\times\vect B)
   &=\vect j\cdot\vect E+\frac{\partial}{\partial t}
      \left(\frac{\varepsilon_0E^2}{2}\right),\\
 \frac1{\mu_0}\vect B\cdot(\nabla\times\vect E)
   &=-\frac{\partial}{\partial t}\left(\frac{B^2}{2\mu_0}\right).
\end{align}
用第二式减去第一式，再利用恒等式
\begin{equation}
  \nabla\cdot(\vect E\times\vect B)
  =\vect B\cdot(\nabla\times\vect E)-
  \vect E\cdot(\nabla\times\vect B),
\end{equation}
可得
\begin{equation}
  \boxed{\frac{\partial u}{\partial t}+\nabla\cdot\vect S=-\vect j\cdot\vect E},
  \qquad
  \boxed{\vect S=\frac{1}{\mu_0}\vect E\times\vect B}.
\end{equation}
这就是电磁场的局部能量守恒式。$\vect S$ 称为 Poynting 矢量，它的方向是能量传播方向，大小表示
单位时间穿过单位垂直面积的能量。

静电场和稳恒磁场可以把能量储存在空间中；当电场与磁场共同组成传播的电磁波时，能量还能够从一处
流向另一处。太阳并不需要把物质一路输送到地球，太阳辐射携带的能量正是由电磁场穿过空间传来的。
因此 Poynting 矢量不是只为配合波动图形而引入的方向箭头，它给出了能量运输的定量描述。对于平面波，
$\vect E\times\vect B$ 指向传播方向，这也反过来帮助判断两场矢量的相对取向。

对有限体积 $V$ 积分，局部守恒式还可以写成
\begin{equation}
  \frac{\dd}{\dd t}\int_Vu\,\dd V
  =-\oint_{\partial V}\vect S\cdot\dd\vect A-
  \int_V\vect j\cdot\vect E\,\dd V.
\end{equation}
左端是体内电磁场能量的增加率；第一个负号项表示由边界流出的电磁能；第二项是场对带电物质做功的
功率。若 $\vect j\cdot\vect E>0$，场把能量交给物质，例如电阻发热；若它为负，外力驱动电荷并向场
输送能量，例如发射天线。$\vect S$ 的单位为 $\si{W/m^2}$，与“单位面积上的功率”这一物理解释一致。

图~\ref{review-energy-balance} 中边界向外的箭头对应面积分，体内同向的电场、电流对应正的做功项。

\poyntingBalanceDiagram

对上面的真空平面波，$E=cB$ 且 $c^2=1/(\mu_0\varepsilon_0)$，所以
\begin{equation}
  u_E=\frac12\varepsilon_0E^2
  =\frac{B^2}{2\mu_0}=u_B.
\end{equation}
电、磁两部分能量在每一时刻都相等。瞬时能流密度为 $S=\varepsilon_0cE^2$，一个周期内的平均光强为
\begin{equation}
  \boxed{\langle S\rangle=\frac12\varepsilon_0cE_0^2
  =\frac{cB_0^2}{2\mu_0}}.
\end{equation}


% EM teaching enrichment 20260927
\paragraph{积分形式中的三个能量去向}
对固定体积 $V$，坡印廷定理写成
\[
 -\oint_{\partial V}\boldsymbol S\cdot\dd\boldsymbol A
 =\frac{\dd}{\dd t}\int_Vu\dd V+\int_V\boldsymbol J\cdot\boldsymbol E\dd V.
\]
左边是从边界净流入的功率，右边分别为内部场能增加率与场对物质的做功率。若是稳恒欧姆导体，场能不变，流入能量最终转化为焦耳热；若是真空充电电容器间隙，没有传导耗散，流入主要增加场能。
对平行板半径 $R$、间距 $d$ 的拟稳充电模型，边缘附近 $B=\mu_0I/(2\pi R)$，侧面流入功率的大小为
$(EB/\mu_0)(2\pi Rd)=IEd=IU$，与电路功率一致。此处忽略边缘修正，只用来核对能量流向和总量。

\subsection{加速电荷的辐射}

静止电荷产生静电场；匀速运动的电荷在固定观察点也可产生随时间变化的场，但在真空中不产生向远区输送净能量的辐射项。电磁作用始终以有限速度传播，不能把“有时间变化”和“有辐射”等同。加速度使场出现按 $1/r$ 衰减的辐射项。在非相对论条件 $v\ll c$ 下，距离电荷很远处的辐射电场大小为
\begin{equation}
  \boxed{E_{\rm rad}(r,t)=\frac{1}{4\pi\varepsilon_0}
  \frac{|q|\,a(t-r/c)\sin\theta}{c^2r}}.
\end{equation}
其中 $t-r/c$ 是推迟时刻，$\theta$ 是观察方向与加速度方向的夹角。这里用 $E_{\rm rad}$ 而不用
$E_r$：辐射场实际上垂直于传播方向，不是沿半径方向的场分量。
$a$ 在大小公式中表示加速度的模；方向应由矢量式判断：
\[
 \vect E_{\rm rad}
 =\frac{q}{4\pi\varepsilon_0c^2r}\,
 \uvec n\times(\uvec n\times\vect a)
 \quad\text{（取推迟时刻）}.
\]
利用三重积，$\uvec n\times(\uvec n\times\vect a)=
\uvec n(\uvec n\cdot\vect a)-\vect a$，恰好是加速度横向分量的反向；
电荷为负时场方向再反转。

辐射场随 $1/r$ 衰减，而 Coulomb 近场随 $1/r^2$ 衰减；所以在足够远处，辐射项反而占主导。远区有
$B_{\rm rad}=E_{\rm rad}/c$，于是单位立体角内的辐射功率为
\begin{equation}
  \frac{\dd P}{\dd\Omega}=r^2S
  =\frac{q^2a^2\sin^2\theta}{16\pi^2\varepsilon_0c^3}.
\end{equation}

图~\ref{review-radiation-pattern} 中角度 $\theta$ 从加速度轴量起，蓝色曲线的径向长度代表功率而不是电场幅值。
图中另画出的 $\vect E_{\rm rad}$ 箭头对应 $q>0$：它垂直于观察方向 $\uvec n$，
沿 $\uvec n\times(\uvec n\times\vect a)$ 指向右下方；若电荷为负，电场方向反转而功率角分布不变。

\radiationPatternDiagram

沿加速度方向 $\theta=0,\pi$ 时不辐射，垂直于加速度时最强。对全空间积分，利用
\begin{equation}
  \int\sin^2\theta\,\dd\Omega
  =2\pi\int_0^\pi\sin^3\theta\,\dd\theta=\frac{8\pi}{3},
\end{equation}
得到 Larmor 公式
\begin{equation}
  \boxed{P=\frac{q^2a^2}{6\pi\varepsilon_0c^3}}.
\end{equation}
它说明辐射功率正比于加速度平方；这里的公式只适用于非相对论点电荷，不能直接套用于
$v$ 接近 $c$ 的情形。

图 10.7 用一组相隔半波长的波前表示振荡电荷向外传递的场。每个蓝箭头都沿波前切向，因而与能流
方向垂直；相邻半波带的箭头反向，对应简谐场的相位翻转。沿加速度轴箭头趋于零，与
$E_{\rm rad}\propto\sin\theta$ 一致。

\radiationWavefrontDiagram

辐射图样常画成绕加速度轴的“甜甜圈”。图中的径向长度表示该方向的相对辐射强弱，并不是电荷真的
喷出一圈物质。把 $E_{\rm rad}$ 代入 $S=\varepsilon_0cE_{\rm rad}^2$，其中 $E_{\rm rad}^2\propto1/r^2$；
再乘球面面积元 $r^2\dd\Omega$，距离因子正好消去。因此无论在多大的球面上统计，只要球面包围波源，
在对应同一推迟时刻的波前上，穿出的总辐射功率都相同。若加速度随时间变化，不同半径处同一观察时刻所收到的是不同发射时刻的能量，不能直接要求瞬时功率相等。这是能量从波源持续向外传播而不是在空间中凭空损失的表现。

推迟时刻 $t-r/c$ 是有限传播速度的直接结果：观察者接收到的不是电荷“此刻”的运动信息，而是场从
电荷传播到观察点之前的状态，正如地球上看到的太阳是约八分钟以前的太阳。Larmor 公式曾是经典
电磁学的重大成功，却也暴露出经典原子模型的根本困难：若电子像行星一样绕原子核转动，它始终有
向心加速度，便应不断辐射并损失机械能，轨道逐渐缩小直至落入原子核。真实原子却稳定存在，这个冲突
不能靠经典电磁学自身解决，后来成为量子理论产生的重要背景之一。

还要分清近场与辐射场。靠近电荷时，$1/r^2$ 的 Coulomb 型项以及与速度有关的感应项可能占主导，
它们储存的能量可在波源附近往返交换；远离波源后，按 $1/r$ 衰减的辐射项成为主导，并携带净能量
离开系统。因此“有时变电场”不等于“观察点处一定只有辐射场”，判断远区辐射必须同时看距离依赖和
能流方向。Larmor 公式计算的是已经离开波源、穿过远处闭合面的总功率。


% EM teaching enrichment 20260927
\paragraph{近场与辐射场的距离检验}
库仑型近场随 $1/r^2$ 衰减，非相对论辐射场随 $1/r$ 衰减。远区能流与场强平方成正比，球面面积又正比于 $r^2$，因此只有 $1/r$ 的场项能在远处留下有限的总辐射功率。
沿加速度轴方向，辐射场的横向投影为零；垂直轴方向最强。图中的 $\sin^2\theta$ 表示每单位立体角的功率分布，不是电场振幅，振幅对应 $\sin\theta$。
若同一非相对论点电荷的加速度大小加倍，拉莫尔功率变为四倍；恒速直线运动的加速度为零，不因观察点处近场随时间改变就产生同样的辐射项。

\subsection{Thomson 散射截面}

考虑沿 $+x$ 传播、沿 $z$ 偏振的电磁波，对初始平均静止的自由电子作弱场、非相对论近似
\begin{equation}
  E_z=E_0\cos(kx-\omega t).
\end{equation}
电子的振动速度幅值约为 $v_0=eE_0/(m\omega)$。要求 $eE_0/(m\omega c)\ll1$，既保证 $v_0\ll c$，也保证位移幅值 $v_0/\omega$ 远小于 $c/\omega$，可把电子所在区域的电场看成均匀。忽略磁力相对于电力的
$v/c$ 阶小量，电子加速度为
\begin{equation}
  a_z=-\frac{eE_0}{m}\cos\omega t.
\end{equation}
把它代入 Larmor 公式并对周期取平均，因 $\langle\cos^2\omega t\rangle=1/2$，有
\begin{equation}
  \langle P_e\rangle=\frac{e^4E_0^2}{12\pi\varepsilon_0m^2c^3}.
\end{equation}
入射波的平均能流密度是 $\langle S_{\rm in}\rangle=\varepsilon_0cE_0^2/2$。散射截面定义为
“总散射功率除以入射能流密度”，于是
\begin{align}
  \sigma_T
  &=\frac{\langle P_e\rangle}{\langle S_{\rm in}\rangle}
  =\frac{e^4}{6\pi\varepsilon_0^2m^2c^4}
  \\
  &=\boxed{\frac{8\pi}{3}\left(\frac{1}{4\pi\varepsilon_0}
  \frac{e^2}{mc^2}\right)^2}
  =\frac{8\pi}{3}r_e^2
  \approx\SI{6.65e-29}{m^2},
\end{align}
其中 $r_e=e^2/(4\pi\varepsilon_0mc^2)$ 称为经典电子半径。截面具有面积量纲，但它不意味着电子是
半径 $\sqrt{\sigma_T/\pi}$ 的实心小球；它衡量的是电子从入射波中截获并再辐射能量的能力。

在这一低频、非相对论经典近似中，$E_0^2$ 在散射功率与入射能流之比中消去，结果也不显含频率，
所以 Thomson 截面与入射波强度、频率均无关。这个结论有明确的适用边界：频率高到必须考虑电子
反冲和量子效应时（光子能量不再满足 $\hbar\omega\ll mc^2$），需要使用更一般的散射理论，不能仍把自由电子当成被均匀电场驱动的经典点电荷。这里还假定电子不受原子束缚；束缚电子在低频、近共振等区域一般不满足这一频率无关结论。
图中的电子电荷为 $-e$，所以瞬时加速度与电场反向；散射强弱取决于 $a^2$，不因这种反向而改变。

散射截面的定义可以想成：若入射能流密度为 $S_{\rm in}$，把总散射功率写成
$P_{\rm sc}=S_{\rm in}\sigma$，那么 $\sigma$ 就是一个等效受光面积。它允许我们不追踪每一束散射波，
而用一个具有面积量纲的量比较相互作用强弱。对一个自由电子，偶极辐射在电子加速度方向上最弱、在
垂直方向上最强；把各方向功率全部积分才得到上面的 $\sigma_T$。经典电子半径 $r_e$ 也不是测得的
电子几何边界，而是由 $e,m,c,\varepsilon_0$ 组合出来的电磁作用长度尺度。

图~\ref{review-thomson} 中向下的加速度与向上的瞬时电场反向；两者沿同一振动轴，散射功率取决于其横向投影。

\thomsonScatteringDiagram


% EM teaching enrichment 20260927
\paragraph{截面的面积单位与物理含义}
散射截面定义为总散射功率除以入射光强，因此单位是面积，但它不是把电子当作半径固定的小硬球得到的几何面积。
经典弱场模型中，电子加速度振幅 $a_0=eE_0/m_e$，平均辐射功率正比于 $a_0^2$，入射光强正比于 $E_0^2$，相除后振幅抵消，得到与入射强度无关的截面。
若频率接近束缚电子的共振，或光子反冲与电子静能不可忽略，自由、非相对论电子的这套推导便不足以使用；它不能无条件推广为所有物质、所有频率的散射规律。

\subsection{超导：零电阻还不够}

金属电阻率一般随温度降低而下降，因为晶格热振动减弱，传导电子受到的散射减少。随着低温技术发展，
人们自然追问电阻会趋向一个非零常数，还是能够真正消失。1908 年，荷兰物理学家 Kamerlingh Onnes（图~\ref{review-onnes-portrait}）
成功液化氦，获得约 $\SI{4.2}{K}$ 的低温条件；1911 年他研究汞的电阻时，发现温度降到这一范围后，
电阻突然跌至实验分辨率以下，超导态由此被发现。

\historyfigure[0.23\linewidth]{onnes-built-in-v1.png}%
{图 10.8：卡末林·昂内斯（Heike Kamerlingh Onnes），低温物理和超导现象的先驱；历史人物肖像。\label{review-onnes-portrait}}

后来人们在多种元素、合金和铜氧化物中发现超导转变。材料进入超导态的温度称为临界温度 $T_c$；
实际超导还受临界磁场和临界电流密度限制，不能只凭 $T_c$ 判断应用性能。从 1911 年的汞到 1973 年前后，
传统材料的 $T_c$ 提升十分缓慢；1986 年贝德诺尔茨（Johannes Georg Bednorz）与缪勒（Karl Alexander M\"uller）发现铜氧化物高温超导体，随后材料的
临界温度迅速越过液氮沸点 $\SI{77}{K}$。下图以分材料族的历史数据展示不同年代的临界温度，反映了
超导材料探索中长期缓慢推进与关键发现后快速跃升并存的过程。

图~\ref{review-superconductivity-timeline} 串起零电阻的发现、迈斯纳效应和新材料温区，图中的年份用于定位各项发现。

\superconductivityTimeline

\superconductivityChartFaithful

1987 年，吴茂昆、朱经武等人的钇钡铜氧材料使转变温度达到约 $90\,$K，超过液氮沸点。随后铋系、铊系和汞系铜氧化物又推动了高温区的探索。早期有过常规超导临界温度约 $30\,$K 的经验估计，但它不是 BCS 理论给出的普适严格上限，不能把后来的高温发现理解为违反一条已证明的上限定理。

零电阻最直接的后果，是闭合超导环中的电流可以长期存在而不发生 Joule 耗散。但若只把超导体看作
电阻率严格为零的普通导体，仍解释不了它最重要的磁学性质。

“电阻测不到”与“电阻严格为零”在实验上并不完全等价，因为任何电压表都有分辨率。超导持续电流
实验提供了另一种证据：在闭合超导环中建立电流后，若电阻哪怕很小，磁能也应按一定时间尺度衰减；
实际观测到电流可以在极长时间内保持，由衰减上限反推的电阻率远低于普通良导体。零电阻意味着
$P_J=I^2R$ 不再产生 Joule 热，却不意味着超导装置可以承受任意电流和磁场。温度、外磁场或电流密度
超过相应临界条件时，材料会回到有耗散的正常态。

超导临界温度的发展也说明实验条件与材料体系同样重要。早期研究依赖液氦低温，能探索的材料范围和
温区都很有限；传统金属与合金的转变温度长期只缓慢提高。铜氧化物高温超导出现后，液氮温区的实验
和应用成为可能，材料的晶体结构、载流子浓度及制备工艺也成为研究重点。图 10.9 将材料名称、年代和实验条件列在各行，用横向长度表示临界温度，避免低温点拥挤。读图时应比较同一时代的数量级与跃迁趋势，
而不宜只盯住某一个最高点；图中高压条件下的数据也不能与常压工程材料直接等同。


% EM teaching enrichment 20260927
\paragraph{持续电流能给出多小的电阻上限}
把环近似为自感 $L$ 与极小电阻 $R$ 的闭合回路，则 $I(t)/I_0=e^{-Rt/L}$。若观察时间 $t$ 内电流相对下降不超过 $\delta$，便有
\[
 R\le-\frac Lt\ln(1-\delta)\simeq\frac{L\delta}{t}\quad(\delta\ll1).
\]
这是一条由测量精度得到的上限，不是由有限精度直接证明一个数学量严格为零。它说明长时间稳定的持续电流为何是灵敏的实验检验，同时也提示应把“测不到变化”与“没有任何变化”在表述上分开。

\subsection{理想导体与迈斯纳效应}

在理想导体内部，$\vect E=0$。Faraday 定律给出
\begin{equation}
  \frac{\dd\Phi_B}{\dd t}=-\oint_C\vect E\cdot\dd\vect l=0.
\end{equation}
因此理想导体只保证磁通不再改变：降温前内部若已有磁场，进入零电阻状态后可以把这部分磁通“冻结”
下来。超导体却表现出更强的性质——在给定温度、外场条件下进入平衡迈斯纳态时，体内磁场被排斥到表面附近，这就是迈斯纳效应。对第二类超导体，需要外场低于相应下临界场；处于混合态时有量子化涡旋进入，不能笼统断言所有超导状态的体内场都为零。

两种操作顺序最能看出差别。若先把无磁场中的样品降至零电阻态、再施加外场，理想导体中的感应电流
会抵消磁通变化，体内可以保持 $B=0$；这个结果尚不能区分理想导体和超导体。若先在正常态样品中
施加磁场，再降温穿过 $T_c$，理想导体模型预言原有磁通应被冻结在体内。1933 年 Meissner 与 Ochsenfeld
测量锡等样品的磁场分布，却发现无论先加磁场还是先降温，只要样品进入超导态且外场低于临界条件，
在平衡迈斯纳态下体内磁场都会被排斥。因而迈斯纳效应是一种平衡态磁学性质，不只是“电阻为零导致磁通不变”。

从磁场边界条件看，体内近似 $B=0$ 必须由表面附近的屏蔽电流来维持。这个电流不是普通电阻回路中
由持续电场推动的电流，而是超导凝聚态对外磁场的无耗散响应。宏观图常把磁感线画成在几何表面突然
转弯，真实情况则是在 London 穿透深度的尺度上连续衰减。样品尺寸远大于穿透深度时，把体内看成
$B=0$ 是很好的竞赛近似；研究薄膜或表面效应时，就必须保留这段有限厚度。

图~\ref{review-meissner} 中左图保留先前的内部磁场，右图外部场线绕开样品；蓝色场线彼此不相交，并满足样品表面无磁通穿入的条件。

\meissnerComparisonDiagram

图中左侧的磁力线穿过材料，表示理想导体保留初始磁通；右侧磁力线绕开超导体，表示迈斯纳排磁。
真实超导体的磁场并非在数学意义上于表面瞬间降为零，而是在称为 London 穿透深度的薄层内衰减；
在本讲的宏观近似中，可写成超导体体内 $B\approx0$。

因此，“完美导电性”和“完全抗磁性”回答的是两个不同问题。前者回答电场能否在体内维持以及电流是否
耗散；后者回答给定温度和外磁场后，平衡态磁通怎样分布。理想导体的最终磁场依赖进入零电阻状态前的
历史，超导体的迈斯纳态则在相同外部条件下趋向同一排磁结果。正是这种与路径无关的排磁性，使超导体
不能被简单等同为电导率无限大的普通金属。

把磁体放在超导体上方时，表面屏蔽电流改变空间磁场分布，可产生向上的磁力。若磁力与重力平衡，
磁体能够悬浮；图~\ref{review-levitation-scene} 中磁体与低温样品间的间隙体现这一状态。实验中常见的稳定悬浮还与第二类超导体的磁通钉扎有关。

图~\ref{review-levitation} 的右侧是磁体的受力图：向上的磁力与向下的重力都作用于该磁体，
平衡时二者大小相等。下方超导圆盘受到的力不属于这个受力分析对象。

\levitationDiagram

\historyfigure[0.32\linewidth]{levitation-built-in-v1.png}%
{超导磁悬浮实验：磁场与屏蔽电流共同产生向上的电磁力。\label{review-levitation-scene}}

\subsection{超导环中的磁通保持}

以下采用宏观厚环、忽略动能电感的近似，并要求全过程不失超、无相位滑移且量子数保持不变。设超导环的几何自感为 $L$，外界穿过环的磁通为 $\Phi_{\rm ext}$，环中持续电流 $I$ 自身贡献磁链 $LI$。
环路电阻为零时，感应电动势满足
\begin{equation}
  L\frac{\dd I}{\dd t}=-\frac{\dd\Phi_{\rm ext}}{\dd t},
\end{equation}
所以
\begin{equation}
  \boxed{\frac{\dd}{\dd t}\bigl(\Phi_{\rm ext}+LI\bigr)=0}.
\end{equation}
先在外磁场中把圆环冷却到超导态，再撤去外场，$\Phi_{\rm ext}$ 减小；环中便出现持续电流，使 $LI$
增加，从而补偿总磁通的变化。

若圆环建立该超导状态时的总磁通为 $\Phi_{\rm ini}$，以后在上述适用条件内缓慢改变外磁场，都有
\begin{equation}
  I=\frac{\Phi_{\rm ini}-\Phi_{\rm ext}}{L}.
\end{equation}
外磁通减小时电流按 Lenz 定律方向增大，外磁通增大时电流反向调整。这里保持的是“外磁通加自感磁通”
的总和，不是说外磁场本身被冻结不变。对于尺度足够大的经典问题，这一磁通守恒关系直接来自零电阻
环路的 Faraday 定律；在微观超导理论中还会出现磁通量子化，那是比本讲宏观模型更进一步的内容。

图~\ref{review-ring-hold} 中外场与自感场均取出纸面为正；逆时针持续电流按右手定则补偿逐渐撤去的外磁通。

\superconductingRingDiagram

这个宏观推导抓住了固定状态分支中总磁链保持的物理图像。$\Phi_{\rm ini}$ 是初始总磁通，不是后文定义的磁通量子 $\Phi_0=h/(2e)$。后面的磁通量子化一节进一步区分总磁通和磁通子，并讨论固定状态与最低能量状态的区别。

% EM teaching enrichment 20260927
\paragraph{撤场以后能量从何而来}
在固定状态分支中，若起初环流为零、外磁通为 $\Phi_i$，随后外磁通降到零，则
$I_f=\Phi_i/L$，最终环内能量为 $\Phi_i^2/(2L)$。
这不是零电阻环无中生有地产生能量；撤场装置与环之间通过互感交换功，必须把外磁场源及其驱动过程一并计入能量账。
零电阻保证没有 $I^2R$ 耗散，却不保证环的磁能不变。若改变几何结构而使 $L$ 改变，也应对 $\Phi_{\rm ext}+LI$ 整体求导，不能继续把 $L$ 提到导数外。

